\documentclass[11pt]{article}
\usepackage[a4paper,margin=1in]{geometry}
\usepackage[T1]{fontenc}
\usepackage[utf8]{inputenc}
\usepackage[english]{babel}
\usepackage{lmodern}
\usepackage{microtype}
\usepackage{amsmath,amssymb,mathtools,array}
\setlength{\parindent}{1.2em}
\setlength{\parskip}{0.35em}
\setlength{\emergencystretch}{8em}
\raggedbottom
\allowdisplaybreaks
\newcommand{\dd}{\,\mathrm d}
\newcommand{\ii}{\mathrm i}
\newcommand{\cosec}{\operatorname{cosec}}
\newcommand{\tg}{\operatorname{tg}}
\newcommand{\arccosop}{\operatorname{arc\,cos}}
\newcommand{\itemtitle}[3]{\clearpage\begin{center}{\Large #1.\\[0.35em] #2}\par\medskip{\normalsize #3}\end{center}}
\newcommand{\datedline}[1]{\begin{flushright}#1\end{flushright}}
\newcommand{\ffone}[6]{F_1\!\left(\begin{array}{ccc}#1&#2&#3\\#4&#5&#6\end{array}\right)}
\begin{document}

\itemtitle{II}{On an Eulerian Integral}{[Journal für reine und angewandte Mathematik, vol.~45, pp.~370--374 (1853).]}

The functions $\Pi$ and $\Gamma$ introduced into analysis by Gauss and Legendre are known to stand in the relation that $\Gamma(a)$ is identical with $\Pi(a-1)$ as long as $a$ has a positive value. For negative $a$, $\Gamma(a)$ is always infinitely large, whereas $\Pi(a-1)$ remains a determinate function and becomes infinite and discontinuous only when $a$ takes one of the values $0,-1,-2$, and so on. The function $\Pi$ is defined as an infinite product, $\Gamma$ as a definite integral. The former definition is unquestionably more comprehensive and gives a deeper insight into the true nature of these functions; nevertheless, for integral calculus it is important to establish a theory of these functions independently, without the aid of those expansions into infinite products and series. This has in fact gradually been accomplished, especially since Dirichlet, in volume 15 of this journal, gave such an elegant proof of Gauss's famous multiplication theorem. In that paper the theorem
\[
\Pi(a-1)\cdot\Pi(-a)=\int_0^\infty \frac{x^{a-1}\dd x}{x+1}=\frac{\pi}{\sin a\pi}
\]
is also used. Very different proofs of it have been given by different mathematicians, but almost all of them proceed through expansions into infinite series. In my dissertation, printed at Easter 1852, \emph{On the Elements of the Theory of the Eulerian Integrals}, the chief ones are collected. There I also added a new route, which stays wholly within the domain of definite integrals, but I was able to give it complete rigor only by presupposing as known the origin of this integral from the multiplication of $\Pi(a-1)$ and $\Pi(-a)$, and the expression for $\frac{\dd\log\Pi(a)}{\dd a}$. In what follows I shall give a proof based on exactly the same idea, but wholly independent of other theories and using only the most general theorems on definite integrals.

First one must recall an auxiliary theorem which will be used several times below. It is well known that
\[
\int\frac{\dd w}{(\alpha w+\beta)(\alpha' w+\beta')}
=\frac{\log\!\left(\frac{\alpha w+\beta}{\alpha' w+\beta'}\right)}{\alpha\beta'-\alpha'\beta},
\]
where the logarithms are hyperbolic. If $\frac{\beta}{\alpha}$ and $\frac{\beta'}{\alpha'}$ are positive quantities, it follows that
\[
\int_0^\infty\frac{\dd w}{(\alpha w+\beta)(\alpha' w+\beta')}
=\frac{\log\frac{\alpha\beta'}{\alpha'\beta}}{\alpha\beta'-\alpha'\beta},
\]
or, if the logarithm is always taken only of the absolute value,
\begin{equation}
\int_0^\infty\frac{\dd w}{(\alpha w+\beta)(\alpha' w+\beta')}
=\frac{\log(\alpha\beta')-\log(\alpha'\beta)}{\alpha\beta'-\alpha'\beta};
\tag{1}
\end{equation}
and this equation is valid even in the case in which $\frac{\alpha}{\beta}=\frac{\alpha'}{\beta'}$, provided the value then assuming the form $\frac00$ is treated according to the rules of differential calculus.

We now turn to the subject proper. First it is easy to see that the given integral
\begin{equation}
\int_0^\infty\frac{x^{a-1}\dd x}{x+1}=A=\varphi(a)
\tag{2}
\end{equation}
has a finite, and indeed positive, value only when $a$ is a positive proper fraction. For if one decomposes the integral into two others with limits $0,1$ and $1,\infty$, and writes $\frac1x$ for $x$ in the latter, one finds
\begin{equation}
A=\int_0^1\frac{x^{a-1}+x^{-a}}{x+1}\dd x.
\tag{3}
\end{equation}
Since throughout the interval of integration $\frac1{x+1}$ lies between the bounds $1$ and $\frac12$, $A$ also lies between the bounds
\[
\int_0^1(x^{a-1}+x^{-a})\dd x
\quad\text{and}\quad
\frac12\int_0^1(x^{a-1}+x^{-a})\dd x.
\]
But this integral has a finite and positive value, namely $\frac1{a(1-a)}$, only when $a$ is a positive proper fraction. The same condition is therefore also necessary for the finiteness of the integral $A$.

Furthermore, equation (3) immediately gives the theorem
\begin{equation}
\varphi(a)=\varphi(1-a).
\tag{4}
\end{equation}
Differentiating this equation with respect to $a$ and then putting $a=\frac12$, one obtains
\begin{equation}
\varphi'\!\left(\frac12\right)=0.
\tag{5}
\end{equation}
Moreover, since $\varphi''(\frac12)$ is positive, as is easily seen, $\varphi(a)$ reaches for $a=\frac12$ a minimum value
\begin{equation}
\varphi\!\left(\frac12\right)
=\int_0^\infty\frac{x^{-1/2}\dd x}{x+1}
=2\int_0^\infty\frac{\dd(\sqrt{x})}{1+(\sqrt{x})^2}=\pi.
\tag{6}
\end{equation}

Let $w$ denote a positive quantity. If in equation (2) one writes $\frac{x}{w}$ instead of $x$, one obtains
\begin{equation}
\int_0^\infty\frac{x^{a-1}\dd x}{x+w}=Aw^{a-1},
\tag{7}
\end{equation}
and if one puts $\frac1w$ for $w$,
\begin{equation}
\int_0^\infty\frac{x^{a-1}\dd x}{xw+1}=Aw^{-a}.
\tag{8}
\end{equation}
Multiplying equation (7) by $\frac{\dd w}{w+1}$, integrating with respect to $w$ between the limits $0$ and $\infty$, and observing that by the auxiliary theorem (1)
\[
\int_0^\infty\frac{\dd w}{(w+1)(w+x)}=\frac{\log x}{x-1},
\]
one obtains
\[
AA=\int_0^\infty\frac{x^{a-1}\dd x}{x-1}\log x.
\]
Now integrate with respect to $a$ between the limits $(1-a)$ and $a$. This gives
\[
\int_{1-a}^{a}AA\dd a
=\int_0^\infty\frac{x^{a-1}-x^{-a}}{x-1}\dd x
=\int_0^\infty\frac{w^{a-1}-w^{-a}}{w-1}\dd w.
\]
Subtracting equation (8) from (7), one easily finds
\[
A\frac{w^{a-1}-w^{-a}}{w-1}
=\int_0^\infty\frac{x^{a-1}(x-1)}{(xw+1)(w+x)}\dd x.
\]
Integrating between the limits $w=0$ and $w=\infty$, and observing that
\[
\int_0^\infty\frac{\dd w}{(xw+1)(w+x)}=\frac{2\log x}{xx-1},
\]
one obtains at once
\[
A\int_{1-a}^{a}AA\dd a
=A\int_0^\infty\frac{w^{a-1}-w^{-a}}{w-1}\dd w
=2\int_0^\infty\frac{x^{a-1}\dd x}{x+1}\log x.
\]
By the property of $A$ that $A=\varphi(a)=\varphi(1-a)$, it follows that
\[
\int_{1-a}^{a}AA\dd a=2\int_{1/2}^{a}AA\dd a.
\]
Moreover,
\[
\int_0^\infty\frac{x^{a-1}\dd x}{x+1}\log x=\frac{\dd A}{\dd a},
\]
and so one finally obtains the equation
\[
A\int_{1/2}^{a}AA\dd a=\frac{\dd A}{\dd a}.
\]
Dividing by $A$ and differentiating with respect to $a$ gives
\[
AA=\frac1A\frac{\dd\dd A}{\dd a^2}-\frac1{AA}\left(\frac{\dd A}{\dd a}\right)^2.
\]
Since the independent variable $a$ does not occur in this equation, introduce $\frac{\dd A}{\dd a}=A'$ as a new variable. This gives
\[
AA=\frac{A'\dd A'}{A\dd A}-\frac{A'A'}{AA},\qquad
A\dd A=\frac{AA\cdot A'\dd A'-A'A'\cdot A\dd A}{A^4},
\]
or
\[
\dd(AA)=\frac{AA\cdot\dd(A'A')-A'A'\cdot\dd(AA)}{(AA)^2},
\]
and the integral of this equation is evidently
\[
AA=\mathrm{Const.}+\frac{A'A'}{AA}
=\mathrm{Const.}+\frac1{AA}\left(\frac{\dd A}{\dd a}\right)^2.
\]
To determine the constant, put $a=\frac12$, for which, by equations (5) and (6), $\frac{\dd A}{\dd a}=0$ and $A=\pi$; hence the value of the constant is $\pi\pi$, and
\[
\dd a=\pm\frac{\dd A}{A\sqrt{AA-\pi\pi}}
=\mp\frac1\pi\frac{\dd(\pi/A)}{\sqrt{1-\frac{\pi\pi}{AA}}}.
\]
If $c$ denotes a constant, integration gives
\[
a+c=\pm\frac1\pi\arccosop\frac\pi A,
\qquad
A=\frac{\pi}{\cos(a+c)\pi}.
\]
To find the constant $c$, put again $a=\frac12$, whence
\[
1=\cos\left(\frac12\pi+c\pi\right)=-\sin c\pi,
\qquad \cos c\pi=0,
\]
\[
\cos(a+c)\pi=\sin a\pi.
\]
Therefore
\[
A=\frac\pi{\sin a\pi},
\]
as was to be proved. In addition, several related integrals follow very easily from this proof; I do not pursue them here.

\datedline{Braunschweig, September 1852.}

\itemtitle{III}{A Theorem from the Theory of Triaxial Coordinate Systems}{[Journal für reine und angewandte Mathematik, vol.~50, pp.~272--275 (1855).]}

If the angles $YOZ$, $ZOX$, $XOY$ of a triaxial coordinate system are denoted by $a,b,c$, then the concave angle $w$ between two arbitrary directions $OM$ and $OM'$ is determined by the following equation:
\begin{equation}
\left\{
\begin{aligned}
&\alpha\alpha'\sin^2a+\beta\beta'\sin^2b+\gamma\gamma'\sin^2c
 +(\beta\gamma'+\gamma\beta')(\cos b\cos c-\cos a)\\
&\quad +(\gamma\alpha'+\alpha\gamma')(\cos c\cos a-\cos b)
 +(\alpha\beta'+\beta\alpha')(\cos a\cos b-\cos c)\\
&\hspace{18em}=D\cos w,
\end{aligned}\right.
\tag{1}
\end{equation}
where $\alpha,\beta,\gamma$ are the cosines of the concave angles $MOX$, $MOY$, $MOZ$, and likewise $\alpha',\beta',\gamma'$ are the cosines of the concave angles $M'OX$, $M'OY$, $M'OZ$, while $D$ has the following meaning:
\begin{equation}
D=1-\cos^2a-\cos^2b-\cos^2c+2\cos a\cos b\cos c.
\tag{2}
\end{equation}
This known theorem includes the other theorem that three such direction cosines, such as $\alpha,\beta,\gamma$, must always satisfy the condition
\begin{equation}
\left\{
\begin{aligned}
&\alpha^2\sin^2a+\beta^2\sin^2b+\gamma^2\sin^2c
 +2\beta\gamma(\cos b\cos c-\cos a)\\
&\quad+2\gamma\alpha(\cos c\cos a-\cos b)+2\alpha\beta(\cos a\cos b-\cos c)=D.
\end{aligned}\right.
\tag{3}
\end{equation}

If the coordinate system is rectangular, equations (1) and (3) become
\[
\alpha\alpha'+\beta\beta'+\gamma\gamma'=\cos w,
\qquad
\alpha^2+\beta^2+\gamma^2=1.
\]
Thus, to express that the three lines $OM,OM',OM''$ then form a second rectangular coordinate system, the following six equations are necessary:
\begin{equation}
\left\{
\begin{aligned}
\alpha^2+\beta^2+\gamma^2&=1, & \alpha'\alpha''+\beta'\beta''+\gamma'\gamma''&=0,\\
\alpha'^2+\beta'^2+\gamma'^2&=1, & \alpha''\alpha+\beta''\beta+\gamma''\gamma&=0,\\
\alpha''^2+\beta''^2+\gamma''^2&=1, & \alpha\alpha'+\beta\beta'+\gamma\gamma'&=0.
\end{aligned}\right.
\tag{4}
\end{equation}
They are also sufficient for this purpose if it is assumed that the first system is rectangular.

The theorem announced in the title consists in the fact that this latter restriction may be omitted, since equations (4) undoubtedly express that both systems must be wholly rectangular. The proof of this remarkable theorem is the subject of the present paper.

First, without further proof, let the known consequences of equations (4) be recorded, namely
\begin{equation}
\left\{
\begin{aligned}
\alpha^2+\alpha'^2+\alpha''^2&=1, & \beta\gamma+\beta'\gamma'+\beta''\gamma''&=0,\\
\beta^2+\beta'^2+\beta''^2&=1, & \gamma\alpha+\gamma'\alpha'+\gamma''\alpha''&=0,\\
\gamma^2+\gamma'^2+\gamma''^2&=1, & \alpha\beta+\alpha'\beta'+\alpha''\beta''&=0,
\end{aligned}\right.
\tag{5}
\end{equation}
and
\begin{equation}
\left\{
\begin{aligned}
\beta'\gamma''-\beta''\gamma'&=\varepsilon\alpha, &
\gamma'\alpha''-\gamma''\alpha'&=\varepsilon\beta, &
\alpha'\beta''-\alpha''\beta'&=\varepsilon\gamma,\\
\beta''\gamma-\beta\gamma''&=\varepsilon\alpha', &
\gamma''\alpha-\gamma\alpha''&=\varepsilon\beta', &
\alpha''\beta-\alpha\beta''&=\varepsilon\gamma',\\
\beta\gamma'-\beta'\gamma&=\varepsilon\alpha'', &
\gamma\alpha'-\gamma'\alpha&=\varepsilon\beta'', &
\alpha\beta'-\alpha'\beta&=\varepsilon\gamma'',
\end{aligned}\right.
\tag{6}
\end{equation}
where, as is known, $\varepsilon^2=1$.

The ternary quadratic form
\[
F=x^2+y^2+z^2+2yz\cos a+2zx\cos b+2xy\cos c
\]
[which, as is known, expresses the square of the distance of an arbitrary point $(xyz)$ from the origin $O$ of the coordinate system $OXYZ$] has as determinant the expression denoted above in (2) by $D$ (the square of the volume of the parallelepiped formed with the three axes $OX=OY=OZ=1$ as edges) and has as adjoint form
\[
\begin{aligned}
F_1={}&x^2\sin^2a+y^2\sin^2b+z^2\sin^2c
+2yz(\cos b\cos c-\cos a)\\
&+2zx(\cos c\cos a-\cos b)+2xy(\cos a\cos b-\cos c).
\end{aligned}
\]
Then the determinant of $F_1$ is, as is known, the square of that of $F$, hence $=DD$, and the adjoint form $F_2$ of $F_1$ is $=DF$.

If one introduces the notation
\[
\begin{aligned}
&xx'\sin^2a+yy'\sin^2b+zz'\sin^2c+(yz'+zy')(\cos b\cos c-\cos a)\\
&\quad +(zx'+xz')(\cos c\cos a-\cos b)+(xy'+yx')(\cos a\cos b-\cos c)\\
&\hspace{13em}\equiv \ffone{x}{y}{z}{x'}{y'}{z'},
\end{aligned}
\]
it is further known from the theory of ternary forms that
\[
\begin{aligned}
&\ffone{x}{y}{z}{x}{y}{z}\,\ffone{x'}{y'}{z'}{x'}{y'}{z'}
-\left[\ffone{x}{y}{z}{x'}{y'}{z'}\right]^2\\
&\quad\equiv F_2\!\left(\begin{array}{ccc}
yz'-zy'&zx'-xz'&xy'-yx'\\
yz'-zy'&zx'-xz'&xy'-yx'
\end{array}\right),
\end{aligned}
\]
and therefore, in the present case,
\begin{equation}
\equiv D\,F\!\left(\begin{array}{ccc}
yz'-zy'&zx'-xz'&xy'-yx'\\
yz'-zy'&zx'-xz'&xy'-yx'
\end{array}\right)
\tag{7}
\end{equation}
holds.

After these preliminaries it is easy to prove the theorem above. Let $OXYZ$ be one coordinate system with angles $a,b,c$, and $OMM'M''$ the other with angles $m,m',m''$. Let $OM$ form with the three axes $OX,OY,OZ$ angles whose cosines are $\alpha,\beta,\gamma$, and similarly for the others. Then the following six equations hold:
\begin{equation}
\left\{
\begin{aligned}
\ffone{\alpha}{\beta}{\gamma}{\alpha}{\beta}{\gamma}&=D, &
\ffone{\alpha'}{\beta'}{\gamma'}{\alpha'}{\beta'}{\gamma'}&=D,\\
\ffone{\alpha''}{\beta''}{\gamma''}{\alpha''}{\beta''}{\gamma''}&=D, &
\ffone{\alpha'}{\beta'}{\gamma'}{\alpha''}{\beta''}{\gamma''}&=D\cos m,\\
\ffone{\alpha''}{\beta''}{\gamma''}{\alpha}{\beta}{\gamma}&=D\cos m', &
\ffone{\alpha}{\beta}{\gamma}{\alpha'}{\beta'}{\gamma'}&=D\cos m'',
\end{aligned}\right.
\tag{8}
\end{equation}
which generally express the relation between any two triaxial coordinate systems. If, however, equations (4), and hence also (5) and (6), hold in addition, then by adding the first three equations in (8) one obtains
\begin{equation}
\sin^2a+\sin^2b+\sin^2c=3D.
\tag{9}
\end{equation}
Moreover, from the theorem contained in (7),
\[
\begin{aligned}
&\ffone{\alpha}{\beta}{\gamma}{\alpha}{\beta}{\gamma}\,
 \ffone{\alpha'}{\beta'}{\gamma'}{\alpha'}{\beta'}{\gamma'}
 -\left[\ffone{\alpha}{\beta}{\gamma}{\alpha'}{\beta'}{\gamma'}\right]^2\\
&=D\,F\!\left(\begin{array}{ccc}
\beta\gamma'-\beta'\gamma&\gamma\alpha'-\gamma'\alpha&\alpha\beta'-\alpha'\beta\\
\beta\gamma'-\beta'\gamma&\gamma\alpha'-\gamma'\alpha&\alpha\beta'-\alpha'\beta
\end{array}\right)\\
&=D\,F\!\left(\begin{array}{ccc}
\varepsilon\alpha''&\varepsilon\beta''&\varepsilon\gamma''\\
\varepsilon\alpha''&\varepsilon\beta''&\varepsilon\gamma''
\end{array}\right)
\end{aligned}
\]
or
\[
DD-DD\cos^2m''
=D(\alpha''^2+\beta''^2+\gamma''^2+2\beta''\gamma''\cos a+2\gamma''\alpha''\cos b+2\alpha''\beta''\cos c),
\]
so
\[
\begin{aligned}
D\sin^2m&=1+2\beta\gamma\cos a+2\gamma\alpha\cos b+2\alpha\beta\cos c,\\
D\sin^2m'&=1+2\beta'\gamma'\cos a+2\gamma'\alpha'\cos b+2\alpha'\beta'\cos c,\\
D\sin^2m''&=1+2\beta''\gamma''\cos a+2\gamma''\alpha''\cos b+2\alpha''\beta''\cos c.
\end{aligned}
\]
Adding gives
\[
D(\sin^2m+\sin^2m'+\sin^2m'')=3.
\]
Comparing this relation with that contained in (9), one obtains
\[
(\sin^2a+\sin^2b+\sin^2c)(\sin^2m+\sin^2m'+\sin^2m'')=3\cdot3,
\]
and hence
\[
\sin^2a=\sin^2b=\sin^2c=\sin^2m=\sin^2m'=\sin^2m''=1;
\]
that is, all six coordinate angles must be right angles.

In this proof it was of course assumed that $D$ is different from zero, that is, that $OX,OY,OZ$ are not contained in one plane.

\datedline{Göttingen, 15 July 1854.}

\itemtitle{IV}{Remarks on a Problem in the Calculus of Probabilities}{[Journal für reine und angewandte Mathematik, vol.~50, pp.~268--271 (1855).]}

In one of the earlier issues of the \emph{Philosophical Magazine} for 1854, G. Boole attacked a solution, given by A. Cayley, of a problem in probability. Although there can be no doubt that the error contained in this attack has already been recognized by others, the following remarks may perhaps not be unwelcome to those interested in this beautiful theory.

The problem is as follows. Given the probability $\alpha$ that a cause $A$ (which can produce a certain event) comes into operation, and the probability $p$ that, if $A$ acts, the event occurs; likewise the probability $\beta$ that a cause $B$ comes into operation, and the probability $q$ that, if $B$ acts, the event occurs: find the probability $u$ of the event, under the assumption that it can be produced by no cause other than $A$ and $B$.

Cayley solves the problem as follows. Let $\lambda$ be the probability that, if $A$ acts, the event is also produced by $A$; let $\mu$ be the probability that, if $B$ acts, the event is also produced by $B$. Then
\[
p=\lambda+(1-\lambda)\mu\beta,
\qquad
q=\mu+(1-\mu)\lambda\alpha.
\]
From these, $\lambda,\mu$ are determined, and the required probability is
\[
u=\lambda\alpha+\mu\beta-\lambda\mu\alpha\beta.
\]
After Boole had found this solution correct under several specializations, he sought to show that in the case $p=1$, $q=0$ it leads to a false result. He says that it is evident that the probability of the event in this case must be $=\alpha$. For if cause $A$ always produces the event, cause $B$ never does, and the occurrence of the event cannot be ascribed to any other cause, then the probability of the event must equal that of the occurrence of cause $A$. Since nothing can of course be objected to in this assertion, and since the solution, as Cayley presents it, gives in this case either $u=1$ or $u=\alpha(1-\beta)$, Boole concludes that the whole solution must be faulty, and gives the final formula of his own solution, with the addition of special restrictions, from which in fact the desired result $u=\alpha$ can be derived for this case.

It is, however, by no means apparent where Cayley has made an error; and indeed his solution is strictly correct, even in the case just mentioned, except for certain restrictions by which it must first be made unambiguous. For one easily finds that $\alpha(1-\beta)$ agrees with $\alpha$, since $\alpha$ can be nothing other than zero. If the possibility of the occurrence of cause $A$ were left open, that is, if $\alpha$ were not zero, then the probability $q$ of the event (under the assumption that cause $B$ occurs) could not possibly vanish entirely, however small $p$ might be, provided only that it is not zero (whereas in the present case $p=1$ was assumed). The proposed problem is therefore nonsensical if $q=0$ while $\alpha$ and $p$ are both assumed to be different from zero. This also follows by merely looking at Cayley's equations. If one observes that $\mu,(1-\mu),\lambda,\alpha$ cannot, by the nature of their meaning, be negative, then from the equation $q=0$ follows both $\mu=0$ and $\lambda\alpha=0$, while the other equation becomes $p=\lambda$. If $p$ is now different from zero (it is not necessary that $p$ be exactly $=1$), then $\alpha$ must also be $0$; and the required probability $u$ must always be $=0$, whether $q$ or $p$, or both, are $=0$, just as one should expect.

But although this objection does not strike the solution above, the solution must still be called incomplete, since the conditions are not given under which the problem has a real meaning, and since it remains to decide which of the two values of $u$ satisfying the above equations is to be chosen. This shall be done here.

The most symmetric procedure is to eliminate $\mu$ from the equations for $q$ and $u$, and likewise $\lambda$ from the equations for $p$ and $u$. This gives
\begin{equation}
u-\beta q=(1-\beta)\lambda\alpha,
\qquad
u-\alpha p=(1-\alpha)\mu\beta,
\tag{1}
\end{equation}
and, substituting these values of $\lambda\alpha$ and $\mu\beta$ into the equation for $u$, one obtains a quadratic equation whose solution gives
\begin{equation}
u=\frac12(1-\alpha\beta+\alpha p+\beta q-\rho),
\tag{2}
\end{equation}
where $\rho$ is the still ambiguous square root of
\[
\rho\rho=(1-\alpha\beta+\alpha p+\beta q)^2-4(1-\beta)\alpha p
-4(1-\alpha)\beta q-4\alpha p\cdot\beta q.
\]
For the problem to be soluble, it is necessary first that $\rho$ be real, and then that $u$, as a probability, be a positive proper fraction. But this is still not enough, and here lies the chief interest of the whole problem. It would still remain meaningless if the auxiliary probabilities $\lambda,\mu$ were not also contained between the limits $0$ and $1$, and it is clear that these latter conditions must at the same time include the former. It remains only to set up the conditions expressing that $\lambda,\mu$ do not lie outside those limits. This is easy, since the values of $\lambda,\mu$ are obtained from equations (1) after substituting into them the expression for $u$ found in (2). In this investigation one arrives at the following equations:
\[
\begin{aligned}
\rho\rho&=(1-2\alpha+\alpha\beta+\alpha p-\beta q)^2
+4\alpha(1-\alpha)(1-\beta)(1-p)\\
&=(1-2\beta+\alpha\beta-\alpha p+\beta q)^2
+4\beta(1-\beta)(1-\alpha)(1-q)\\
&=(1-\alpha\beta+\alpha p-\beta q)^2-4\alpha(1-\beta)(p-\beta q)\\
&=(1-\alpha\beta-\alpha p+\beta q)^2-4\beta(1-\alpha)(q-\alpha p).
\end{aligned}
\]
From the first two forms for $\rho\rho$ it follows that no special condition is needed for the reality of $\rho$. But if one places the forms in connection with the requirements for $\lambda,\mu$, it follows that in the expression (2) for $u$ the positive square root must always be taken for $\rho$. Finally, comparing the last two forms for $\rho\rho$ with the requirements for $\lambda$ and $\mu$, one obtains, as the only necessary and also wholly sufficient conditions, that the two differences
\begin{equation}
p-\beta q\quad\text{and}\quad q-\alpha p
\tag{3}
\end{equation}
may not be negative.

Thus if, in order to give an example of the problem, one has assumed for $\alpha,\beta,p,q$ four arbitrary numbers between the bounds $0$ and $1$, one must first examine whether they satisfy the two conditions in (3). Incidentally, in such an arbitrary choice one can just as often encounter a nonsensical example as a suitable one; for the value of the fourfold integral
\[
\int\!\int\!\int\!\int \dd\alpha\,\dd\beta\,\dd p\,\dd q
\]
is $=1$ if the integrations are extended over all values of the variables between $0$ and $1$, but $=\frac12$ if one excludes those values which do not satisfy the conditions (3). This can also be seen easily by geometric considerations.

In the case examined by Boole, $q=0$, the conditions (3) reduce to $\alpha p=0$; then $\rho=1-\alpha\beta$, and hence $u=0$, in complete agreement with the above results. The case $\alpha=0$ is also of interest. Then $\rho=1-\beta q$, and hence $u=\beta q$ is plainly the correct result. Here $u$ is of course independent of $q$, and nevertheless the condition $p-\beta q\ge0$ remains in full force. Although the value of $p$ carries no weight at all for determining $u$, it would still be nonsensical to assume the probability $p$ of the event under the assumption that cause $A$ comes into operation to be smaller than the probability $\beta q$, even though that assumption is factually forbidden by setting $\alpha=0$. With this remark, which I believe is suited to cast a striking light on the peculiar character of this kind of problem, I shall break off my considerations.

\datedline{Göttingen, 22 July 1854.}

\end{document}
